\section{Small-level projection and identification of the base theta map}
\label{sec:small-level-rank}

\status{SLR1--SLR9 reconstruct the original level projection and the finite
rank argument.  The projection is onto on every supplied
\texttt{Bundle}, with an explicit right inverse and no additional size
threshold.  For the original selected row family, the remaining sufficient
input is its base row realization together with its theta comparison at
its base tuple.  From these inputs we prove the literal
\texttt{SmallLevelRankEq} conclusion, including the quantifiers of
\texttt{Closed}.  We do not claim to have constructed those actual row
and theta inputs.  This replaces the former opaque projection question
by an explicit, narrower producer boundary.}

The target is an equality of complex dimensions for the original record:
\[
 \exists x_0\in\mathbb R\ \forall x\ge x_0\ \forall i\in\mathcal F.I_x\
 \ \forall j\in\{1,\ldots,(\mathcal F.D_{x,i}).R.m\},\qquad
       dZ_0(j)=D_0.                                      \tag{SLR.1}
\]
Here $dZ_0(j)=\dim_{\mathbb C}Z_j(\mu_0)$ and
$Z_j(\mu)=\operatorname{im}\theta_{\mu,j}^{*}$, with
$\theta_{\mu,j}:M_t(\mu,j)\to M_j$ the actual map stored in
$D.\texttt{what}_2$.  The superscript $*$ in this expression means
Hilbert-space adjoint.  Matrix conjugate transpose will also be denoted
by $*$; their roles will always be specified.  The separate source
symbol \texttt{lZ0} denotes a coefficient matrix, not the complex
space $Z_j$.

All source references below use the pinned revision.  The accompanying
\texttt{provenance.json} gives exact paths, ranges and hashes.  The principal target references are S01--S03, S27
and S37--S38.  Existing CR, CS and TD arguments remain conditional on
their stated actual data; an abstract rank equality does not supply
that data.

\subsection{The coefficient projection has an explicit right inverse}

% BEGIN SLR1
\paragraph{SLR1: the actual level and good spaces are nested.}
\label{claim:SLR1}
Let $K$ be the coefficient number field, let $\sigma:K\hookrightarrow
\mathbb C$ be its selected embedding, and assume the coefficient
involution satisfies $\sigma(\bar a)=\overline{\sigma(a)}$.  Let a
finite group $G$ act on a finite-dimensional $K$-space.  For a subgroup
$P$ and a character $\psi:P\to K^{\times}$, set
\[
 e=\frac1{|P|}\sum_{h\in P}\psi(h^{-1})[h],\qquad
 C=\sum_{\chi\in X}\frac{\chi(1)}{|G|}
                  \sum_{g\in G}\chi(g^{-1})[g],
 \qquad a=Ce.                                           \tag{SLR.2}
\]
The finite set $X$ consists of characters of absolutely irreducible
representations, as in the source's \texttt{IsAbsIrrChar}.  Each such
character is constant on conjugacy classes: over the algebraic closure
it is a matrix trace, and conjugation replaces its representing matrix
by a similar matrix.  Injectivity of $K\to\overline K$ brings this
equality back to $K$.  Therefore $C$ is central in $K[G]$.

For $u\in P$, the coefficient of $[u]$ in $e^2$ is
\[
 \frac1{|P|^2}\sum_{h\in P}
          \psi(h^{-1})\psi((h^{-1}u)^{-1})
       =\frac1{|P|}\psi(u^{-1});
\]
there are no coefficients outside $P$.  Thus $e^2=e$, and centrality
gives $e(Ce)=Ce$.  Applying the group action shows
\[
           \operatorname{im}\rho(a)\subseteq
                         \operatorname{im}\rho(e).       \tag{SLR.3}
\]
This argument uses neither character orthogonality nor any equality of
integral lattices.  It also works when either space is zero.
\emph{Source: S04, S05, S07.}
% END SLR1

% BEGIN SLR2
\paragraph{SLR2: the matrices defining \texttt{lPrj} are onto.}
\label{claim:SLR2}
Use the original coefficient basis of the finite form space, of
dimension $n$, and let $H$ be its Hermitian matrix.  At $\sigma$ it is
positive definite.  Let $L\in M_{n,\ell}(K)$ have as columns the source's
chosen basis of $\operatorname{im}\rho(e)$, and let
$V\in M_{n,v}(K)$ have as columns its chosen basis of
$\operatorname{im}\rho(a)$.  For the original subspace $Z$ of forms,
let $U\in M_{v,d}(K)$ be the chosen basis matrix of the inverse image
of $Z$ under $V$.  These three matrices have full column rank by their
definitions as subspace basis matrices.  Define
\[
 \begin{aligned}
 L^+&=(L^*HL)^{-1}L^*H,&
 V^+&=(V^*HV)^{-1}V^*H,\\
 H_V&=|P|^{-1}V^*HV,&
 Q&=(U^*H_VU)^{-1}U^*H_V,\\
 E&=L^+VU,& \Pi&=QV^+L.
 \end{aligned}                                             \tag{SLR.4}
\]
These are precisely \texttt{lvlPr}, \texttt{vPr}, \texttt{HV},
\texttt{Pq}, \texttt{lE} and \texttt{lPr} of the source.

For clarity, all inverses in (SLR.4) exist over $K$.  If a matrix $A$
has full column rank, some maximal minor is nonzero in $K$, hence also
at $\sigma$.  Thus $\sigma(A)$ is injective.  Positivity of $\sigma(H)$
makes $\sigma(A)^*\sigma(H)\sigma(A)$ positive definite, hence
nonsingular.  Its determinant is the image of
$\det(A^*HA)$, so this determinant is nonzero in $K$.  The same argument
applies to $H_V$, since $|P|^{-1}>0$ at $\sigma$.  The empty matrix
case has determinant $1$ and causes no exception.

Multiplying out gives $L^+L=1$, $V^+V=1$ and $QU=1$.
By (SLR.3), every column of $V$ lies in the image of $L$;
consequently $LL^+V=V$.  Therefore
\[
             LE=VU,\qquad
             \Pi E=QV^+LL^+VU=QV^+VU=QU=1.              \tag{SLR.5}
\]
The map $y\mapsto Ey$ is an explicit right inverse to $x\mapsto\Pi x$.
This proves surjectivity, not just equality of source and target ranks.

There is also an exact orthogonal projection identity.  Put
$H_L=|P|^{-1}L^*HL$.  From $LE=VU$ we get
\[
 E^*H_LE=U^*H_VU,
 \qquad
 E^*H_L=|P|^{-1}U^*V^*HL.
\]
After cancelling the nonzero factor $|P|^{-1}$ in (SLR.4), this yields
\[
             \Pi=(E^*H_LE)^{-1}E^*H_L.                  \tag{SLR.6}
\]
In particular $E^*H_L(x-E\Pi x)=0$ for every $x$.  Thus $E\Pi$ is
projection onto the columns of $E$ along their $H_L$-orthogonal
complement.  This verifies the normalization by $|P|$ as well as the
right inverse.
\emph{Source: S04--S08.}
% END SLR2

% BEGIN SLR3
\paragraph{SLR3: the precise lattice consequence and its limit.}
\label{claim:SLR3}
For an $\mathcal O_K$-submodule $N\subset K^n$, write
$\operatorname{cut}_A(N)=\{u:Au\in N\}$.  Let $N$ be the actual
full coefficient lattice of forms, and let $N'\subseteq N$ be the
source's smaller lattice.  With (SLR.4), the source's level and lift
lattices are
\[
 N_L=\operatorname{cut}_L(N),\qquad
 T=\Pi N_L,\qquad
 N_Z=\operatorname{cut}_{VU}(N),\qquad
 N'_Z=\operatorname{cut}_{VU}(N').                       \tag{SLR.7}
\]
The exact identity $LE=VU$ gives
$\operatorname{cut}_E(N_L)=N_Z$.  If $z\in N_Z$, then $Ez\in N_L$
and $\Pi Ez=z$.  Hence
\[
                         N'_Z\subseteq N_Z\subseteq T.   \tag{SLR.8}
\]

Here $N_L$, $N_Z$, $N'_Z$ (when $N'$ is full), and $T$ are full
lattices in their respective coefficient spaces.  One elementary
verification is as follows.  If $A:K^d\hookrightarrow K^n$ is
injective, choose a $K$-linear left inverse $B$.  Then
$\operatorname{cut}_A(N)\subseteq B(N)$; the latter is finitely
generated over $\mathcal O_K$, and its submodule is finitely generated
because $\mathcal O_K$ is Noetherian.  A common nonzero denominator
sends each vector $Ae_r$ into $N$, so the cut contains a nonzero scalar
multiple of every coordinate vector.  It spans $K^d$.
Images of finite generating sets generate $\Pi N_L$, and
$K(\Pi N_L)=\Pi(KN_L)=K^d$ by (SLR.5).  These observations prove the
assertions.

The surjectivity in SLR2 is over $K$.  We have not proved that $E$
preserves any independently specified integral lattice, or that
$N_Z=T$.  The map $N_L\to T$ is onto because $T$ is its image by
definition.  No stronger integral assertion enters (SLR.1).
\emph{Source: S05, S08.}
% END SLR3

\subsection{The finite type decomposition and the base dimension}

% BEGIN SLR4
\paragraph{SLR4: the actual types give the needed direct sums.}
\label{claim:SLR4}
We give the finite-dimensional argument because the source's type
spaces are sums of joint \emph{generalized} eigenspaces.  No
semisimplicity assumption is necessary.

Let $T_\alpha$ be a commuting family of endomorphisms of a finite
$K$-space $F$, each split over $K$.  Their generated unital algebra
is a finite-dimensional subspace of $\operatorname{End}_K(F)$.
Choose a basis of that algebra made from finitely many products of
operators; the finitely many operators occurring in those products
already generate the entire algebra.  For one such operator, its
minimal polynomial is a product of coprime powers $(X-a)^r$.
Bezout's identity, applied to these factors, supplies polynomial
idempotents that sum to $1$ and project onto their kernels.  All the
other operators preserve those kernels because they commute.  Applying
this construction successively to the finite generating family gives
\[
          F=\bigoplus_{\lambda\in\Lambda}F_\lambda,
          \qquad |\Lambda|<\infty,                      \tag{SLR.9}
\]
with polynomial projections $p_\lambda$ in the operators.
On a resulting block each generator is its scalar eigenvalue plus a
nilpotent operator.  The nilpotent parts commute; their generated ideal
is nilpotent, since a sufficiently long monomial contains a power at
least the nilpotence exponent of one generator.  It follows that every
operator of the original family has a single generalized eigenvalue
on that block.  Blocks with distinct scalar data are disjoint.  Thus
(SLR.9) is precisely the decomposition by the source's intersections
of maximal generalized eigenspaces.  If $F=0$, it is the empty direct
sum.

Every subspace stable under the operators is stable under the
$p_\lambda$.  If the operators commute with the $G$-action, these
projections also commute with $\rho(Ce)$.  Consequently
$Z\cap\operatorname{im}\rho(Ce)$ is the direct sum of its
intersections with $F_\lambda$, whenever $Z$ is stable.
Pulling back under the injective coefficient map $VU$ preserves this
direct-sum decomposition.

We now check the source's labels.  A type $p$ consists of residual
polynomials $p_1(P_0)$ and inert eigenvalues $p_2(I_0,y)$.  A joint
system $\lambda$ has that type if, at each actual Hecke place $P$,
\[
 H_P(\lambda;X)=p_1(\operatorname{pl}P;X)B_P(X),\qquad
 B_P(X)=\prod_{\nu\in\texttt{bl}}\prod_{c=0}^{t-1}
                      (1-\nu(P)q_P^cX),                 \tag{SLR.10}
\]
and its inert eigenvalues are the specified $p_2$ values.
The polynomial $B_P$ has constant term $1$, so it is nonzero.
Hence the residual polynomial, when it exists, is unique by cancellation
in $K[X]$.  The labels are normalized to $1$ at unused Hecke places
and to $0$ at unused inert places.  The place maps are injective, so
these assignments are well defined.  Two distinct normalized labels
therefore differ at an actual place, and cannot label the same
$\lambda$.  Since there are only finitely many occurring $\lambda$,
there are only finitely many occurring labels.  Thus the finite type
set \texttt{ixFin3} really is that set; its conditional definition does
not select its fallback empty set here.

At the base tuple, $Z$ is contained in the prescribed block sum by
\texttt{Zs.InBlock}.  Each nonzero primary component in
$Z\cap\operatorname{im}\rho(Ce)$ therefore satisfies the divisibility
in (SLR.10), and gives an occurring normalized type.  It also belongs
to the full good space used to define \texttt{ixFin3}: the image of
$\rho(Ce)$ is contained in the image of $\rho(C)$, and every
$p_\lambda$ preserves both images.  The nonzero component witnesses
that its type occurs there.  The preceding projections consequently
give
\[
        K^{dQ(t_0)}=\bigoplus_{p\in I}Q_p(t_0).           \tag{SLR.11}
\]
For the lift spaces, distinct transformed labels are different at an
actual lift place by \texttt{TypesSep3}.  The same cancellation and
primary-space argument prove independence of their inverse-image
spaces $Z_{k,p}(t_0)$.  Their spanning is exactly the \texttt{types}
field of the supplied base \texttt{IndefRows} record.  Thus
\[
        K^{dZ(k,t_0)}=\bigoplus_{p\in I}Z_{k,p}(t_0).      \tag{SLR.12}
\]
All hypotheses just used are fields of \texttt{RowsHold0\_3}: stable
and block-contained subspaces, some \texttt{fam3.Ok3 s}, the common
base \texttt{IndefRows}, and type separation.  We have proved the
three pieces of \texttt{WF} needed for dimension transport; its metric
and lattice clauses need not be imported into this argument.
\emph{Source: S09--S14.}
% END SLR4

% BEGIN SLR5
\paragraph{SLR5: the rational lift rank equals the original $D_0$.}
\label{claim:SLR5}
Fix the original selected row family $P$, and a base row witness
$R_0$ with $R_0.\texttt{toIndefData}=P.S_0$ as in
\texttt{RowsHold0\_3}.  Write $D_r=P.S_0.D$ and
$d_k=P.S_0.\texttt{dW}(k)$.  The \texttt{l1c} field gives bijections
\[
 \iota_0:K^{dQ(t_0)}\longrightarrow K^{D_r},\qquad
 \vartheta_{k,0}:K^{dZ(k,t_0)}\longrightarrow K^{d_k}.
                                                               \tag{SLR.13}
\]
The \texttt{l1a} field says these maps carry the summands in
(SLR.11)--(SLR.12) onto $E_p$ and $W_{k,p}$, respectively.
Therefore these latter spaces are direct sums of their indicated
summands.  The \texttt{l1d} field says
$\dim_KW_{k,p}=\dim_KE_p$ for every $p$.  Taking a basis of each
summand and concatenating these bases proves
\[
 d_k=\sum_{p\in I}\dim_KW_{k,p}
       =\sum_{p\in I}\dim_KE_p=D_r.                     \tag{SLR.14}
\]
There is no exchange of an infinite sum or choice of unrelated type
sets in this equality.

The equality $D_r=D_0$ follows directly from data already in the
original \texttt{RowSideFamily}.  Its $P.I.\texttt{eRR0}$ is an
isomorphism
\[
                 K^{D_r}\simeq_K (\texttt{ex.p0.s.RR}).V.
\]
The latter rational space has a specified basis $b_a$, indexed by the
finite set $A_0$.  The field \texttt{ex.p0.s.jEE\_basis} states that
the images $j_{EE}(b_a)$ form a complex basis of the original
$EE_0$.  Finally, \texttt{onb0} is a complex orthonormal basis of
that same $EE_0$, indexed by $\{1,\ldots,D_0\}$.  Counting these
bases gives
\[
          D_r=\dim_K(\texttt{ex.p0.s.RR}).V
              =|A_0|=\dim_{\mathbb C}EE_0=D_0.           \tag{SLR.15}
\]
This uses no isometry between newly chosen spaces.  In particular the
separate \texttt{RowsHoldLam3} predicate is unnecessary for this
numerical comparison.

The \texttt{li} and \texttt{span} fields of
$P.\texttt{tied.div.th}$ assert that the vectors
$j_{W,k}(e_r)$ form a complex basis of the original $M_{k+1}$,
where $r=1,\ldots,d_k$.  Thus (SLR.14)--(SLR.15) imply
\[
                      \dim_{\mathbb C}M_{k+1}=D_0.       \tag{SLR.16}
\]
\emph{Source: S13--S19.}
% END SLR5

\subsection{The exact theta comparison needed for the rank conclusion}

% BEGIN SLR6
\paragraph{SLR6: a base-tuple square makes the actual theta map onto.}
\label{claim:SLR6}
For this paragraph the additional input is the following equality on
the same selected $P$, at its single base tuple $t_0=P.A_3.A.\mu_0$:
\[
 \theta_{P.I.\mathrm{tup}(t_0),\,k+1}
       \bigl(j_{Mt}(e_{Mt,k,t_0}(x))\bigr)
     =j_{W,k}\bigl(\vartheta_{k,0}(\Pi_{k,t_0}x)\bigr)
       \quad\text{for every }k,x.                       \tag{SLR.17}
\]
Here $\Pi_{k,t_0}$ is the original \texttt{lPrj} matrix,
$e_{Mt,k,t_0}=P.I.\texttt{eMt}(k,t_0)$, and the maps $j_{Mt}$ and
$j_{W,k}$ are the actual embeddings in the original record.  Equation
(SLR.17) is the base-tuple restriction of \texttt{thTheta}, not a
consequence of the basis assertions \texttt{li/span} alone.

For each coordinate vector $e_r\in K^{d_k}$ choose its unique
preimage $z_r=\vartheta_{k,0}^{-1}e_r$ under (SLR.13), and put
$x_r=E_{k,t_0}z_r$.  By SLR2,
$\Pi_{k,t_0}x_r=z_r$.  Substitution in (SLR.17) shows that the range
of the actual theta map contains every $j_{W,k}(e_r)$.  These vectors
are a complex basis by SLR5; its range is a complex subspace, so the
map is onto.  In fact (SLR.17) only needs to be checked at this finite
list of $x_r$ for each $k$.  We retain its all-$x$ base-tuple version
as a convenient sufficient contract.

For completeness, if a linear map $A:V\to W$ between finite
complex inner-product spaces is onto and $A^*w=0$, choose $v$ with
$Av=w$.  The adjoint identity gives
$\langle w,w\rangle=\langle Av,w\rangle
=\langle v,A^*w\rangle=0$, hence $w=0$.
Thus $A^*:W\to V$ is injective and its image has dimension
$\dim_{\mathbb C}W$.  Apply this to the actual theta map and use
(SLR.16).  Since $P.\texttt{tied.div.tup.mu0}$ identifies
$P.I.\mathrm{tup}(t_0)$ with the original $\mu_0$, we obtain
\[
       \dim_{\mathbb C}Z_{k+1}(\mu_0)=D_0.               \tag{SLR.18}
\]
If there are no lifts, the requested family of equalities is empty.
Otherwise every $1\le j\le m$ is uniquely $k+1$ with
$k=j-1\in\{0,\ldots,m-1\}$, so (SLR.18) covers exactly the source's
\texttt{SizeData}\textsubscript{2}\texttt{.lifts}.
\emph{Source: S01--S03, S17--S22.}
% END SLR6

% BEGIN SLR7
\paragraph{SLR7: what is still required to identify the actual square.}
\label{claim:SLR7}
The source separates \texttt{DivDef3}, already inside $P.\texttt{tied}$,
from \texttt{DivRows3}.  Its \texttt{ThDef3} contains linear
independence, complex spanning and a lattice tie; it contains no
assertion (SLR.17).  The latter is a field of \texttt{DivRows3}, and
its used-tuple version is a field of \texttt{DivRows3G}.  Since
\[
 \texttt{Used0}(D,\mu)\ \Longleftrightarrow\
       \mu=\mu_0\ \text{or}\ \mu=\texttt{muAt}(w)\text{ for some }w,
                                                               \tag{SLR.19}
\]
the base tuple is used by its first alternative.  Therefore a supplied
\texttt{DivRows3G} for this $P$ proves (SLR.17) directly.  There is no
need to prove \texttt{TupAllUsed}, or to upgrade every other tuple to
\texttt{DivRows3}, for the argument above.

A more detailed reading of the actual theta proof reveals a narrower
analytic boundary.  At an actual realization let $\Theta$ be its
analytic theta map on level forms and let $\epsilon_W:M_{k+1}\to H_W$
be the stored injective reading into the analytic model.  Write
$J(x)$ for the actual level form represented by the rational
coordinates $x$, and $g(y)$ for the actual good form with coordinates
$y$.  The needed comparisons are
\[
 \begin{split}
 \epsilon_W\theta\bigl(j_{Mt}(e_{Mt}x)\bigr)&=\Theta J(x),\\
 \epsilon_W j_W(y)&=g(y),\\
 \Theta J(Ey)&=g(\vartheta y),\\
 E^*H_Lx=0&\Longrightarrow\Theta J(x)=0.
 \end{split}                                             \tag{SLR.20}
\]
These formulas use the same realization, coefficient embedding and
level character as the matrices.  The source's phase is $1$ by its
\texttt{Pins.phs} condition; its possible conjugate scalar model is
transported back before the equalities in (SLR.20) are compared.

Once (SLR.20) is available the remaining implication is elementary.
By (SLR.6), $x-E\Pi x$ is orthogonal to every column of $E$, so the
last line of (SLR.20), together with additivity, gives
$\Theta J(x)=\Theta J(E\Pi x)=g(\vartheta\Pi x)$.
The first two lines and injectivity of $\epsilon_W$ then give
(SLR.17).  This proves the projection-factorization step completely.

We do not label all of (SLR.20) elementary.  In the pinned proof the
first line is the Hilbert-model/lowest-$K$-type reading, whose proof
uses restriction to lowest $K$-types and uniqueness there.  The last
line uses the actual rational adjoint-image and level compatibility.
The third line uses the actual rational theta matrix.  TD proves useful
finite and rational consequences from explicit analytic inputs, but
does not by itself supply all these identifications on the selected
record.  This is a precise remaining human reconstruction obligation.
The proof of \texttt{thTheta} is obtained before the extra pairing
and column fields of full \texttt{DivRows3}; those extra fields are
not hypotheses of the rank implication established here.
\emph{Source: S19--S23, S30--S34.}
% END SLR7

\subsection{Specialization to the selected family and closure}

% BEGIN SLR8
\paragraph{SLR8: the projection certificate is pointwise on the literal bundle.}
\label{claim:SLR8}
Let $b$ be the original \texttt{Bundle} at $(D,O,c)$ and put $P=b.P$.
Its column-compatibility field $b.C.\texttt{ok1}$ includes
$P.A_3.\texttt{fam3.Ok3}(1)$.  Consequently the lift block at any
$k,t$ has: the stated coefficient involution at its selected embedding,
the irreducible character data of SLR1, and the positive Hermitian
matrix used in SLR2.  The exact lift level is
\[
 \begin{aligned}
 B&=P.A_3.A.\mathrm{fam}.L(k,t),\\
 \ell&=P.A_3.A.\mathrm{fam}.\mathrm{lvlA}(P.P_0,P.\psi_0)
          =(X,P.P_0,P.\psi_0^{-1}),\\
 Z&=P.Z_s.Z_l(k,t).
 \end{aligned}                                           \tag{SLR.21}
\]
By definition \texttt{lZ0}$=E$ and \texttt{lPrj}$=\Pi$ for these
objects.  Notice the inverse character on the lift level; replacing
it by the base-level character would change the objects.
SLR2 proves the actual projection surjective, in particular at $t_0$.
The appearance of \texttt{FormsTie} in one source wrapper is not an
extra need for this proof: that wrapper uses only the family
\texttt{Ok} data in its proof body.

A raw bundle selector $\sigma$ gives
$\operatorname{selP}(\mathcal F,h_c,x,h_x,h,i).1
=(\sigma\mathcal F h_c x h_x i).P$ and likewise obtains the columns
and realization from the corresponding $C$ and $B$ fields of this
same bundle.  Its selection threshold is $\mathcal F.x_1$.
Thus all occurrences of $P$, $\Pi$, $E$, $j_W$, and $e_{Mt}$ above
refer to the same literal selection.  The realization used by
\texttt{selIOf} is $(\sigma\mathcal F h_c x h_x i).B.Q$.
There is no second existential choice to reconcile.

For the original \texttt{H\_zL} predicate choose its additional
threshold $x_r=0$.  Choose $\varepsilon_1=1$ and all four components
of its \texttt{Closed} bound $K$ equal to zero.  For every allowed
$\varepsilon,\mathcal F,CH$ the pointwise argument just given proves
the required surjectivity; the largeness premise is unused.  This is
the exact closure of the projection obligation.  No period
admissibility input or new uniform arithmetic bound is required.
\emph{Source: S06, S07, S24--S29, S35--S36.}
% END SLR8

% BEGIN SLR9
\paragraph{SLR9: a narrower sufficient producer cut closes the original rank predicate.}
\label{claim:SLR9}
Fix a class $\mathcal A$ and the original compatible selectors.  It
suffices that, for each allowed family $\mathcal F\in\mathcal A$,
there is a real number $a_{\mathcal F}$ such that every selected
$P$ at $x\ge a_{\mathcal F}$ with $x>\mathcal F.x_1$ has both:
\begin{enumerate}
\item \texttt{RowsHold0\_3} for this $P$;
\item the base-tuple equality (SLR.17) for this same $P$.
\end{enumerate}
Here $a_{\mathcal F}$ is also taken at least the selector threshold;
if the two inputs initially have different eventual thresholds,
replace them by their maximum.  This is one common-family contract,
not two assertions that possibly use different row realizations.
The source obtains these inputs from the stronger hypotheses
\texttt{PrimAll} and pointwise \texttt{DivRows3G}, with their original
common selector.  SLR7 proves the restriction of the latter to the
base tuple.  On the row side, the primitive-row construction that turns
\texttt{PrimAll} into an \texttt{IndefRows} witness remains part of
the producer boundary here.  Once that witness has been supplied, its
transport uses the realization's actual equalities of $Z_s$ and $S_0$:
substitute the first equality in the type of the witness, then the
second in its underlying data.  The resulting witness has exactly
$P.S_0$ as underlying data.  No isomorphism of unrelated abstract row
spaces is substituted for these equalities.

Choose
\[
              x_0=\max\{a_{\mathcal F},\mathcal F.x_1+1\}.
                                                               \tag{SLR.22}
\]
If $x\ge x_0$, the strict condition $x>\mathcal F.x_1$ holds.
Every index $i\in\mathcal F.I_x$ has the same compatible objects
as in SLR8, the two producer inputs hold, and SLR4--SLR6 give
(SLR.18) for every $j\in\{1,\ldots,m\}$.  This proves (SLR.1) for
the original family, with one threshold for all $i,j$.

It also proves the source's
$\texttt{Closed}(D3\ 3\ \mathcal A)\,
(\mathcal F,CH\mapsto\texttt{SmallLevelRankEq}\,\mathcal F)$.
Indeed, choose $\varepsilon_1=1$ and the zero bound $K$ before
$\mathcal F,CH$ are supplied.  Membership in $D3\ 3\ \mathcal A$
forgets to membership in $\mathcal A$; apply the preceding argument.
The only family-dependent threshold is $x_0$ inside the existential
quantifier of \texttt{SmallLevelRankEq}, exactly where the definition
permits it.  No threshold is moved in front of a family quantifier.

\emph{Audit consequence.}  The logical implication from this common
producer cut to the original \texttt{GLOBAL\_SMALL\_RANK} target is
now completely reconstructed.  Its projection premise is discharged
by SLR1--SLR3 and SLR8, and its finite rank and threshold arguments by
SLR4--SLR6 and this paragraph.  The actual base row and theta producer
remains open in the human reconstruction.  Its arithmetic and analytic
construction is shared with the actual row/period program; such a
shared work relationship is not itself a logical proof of either
input.  The construction of the bundle/family also remains an input.
In particular, this chapter supplies the rank input to
\texttt{topDesWin} conditionally on the displayed actual producer;
it does not close every other input of that theorem.
\emph{Source: S01, S15, S20--S29, S35--S39.}
% END SLR9
